% assets/w05-los-standard.svg — LOS 유도의 기하. Fossen 의 표준 그림 배치.
%
%   bash _tools/tikz2svg.sh assets/src/w05-los-standard.tex assets/w05-los-standard.svg
%
% 배치와 기호는 Fossen 의 TTK 4190 "Guidance, Navigation and Control of
% Vehicles" 강의노트, 그리고 Handbook of Marine Craft Hydrodynamics and Motion
% Control (2021, 2nd ed.) §12.3 을 따른다. 대학원 과목 "센서신호처리 및 융합"
% W05 의 같은 그림을 **배치와 좌표를 그대로 두고 글자만 한글로** 옮긴 것이다
% (지시 2026-10-20).
%
%   p_i^n, p_{i+1}^n   웨이포인트 (NED). 굵은 소문자에 위첨자 n
%   pi_p               경로 접선각
%   x_e^p, y_e^p       경로좌표계 {p} 에서의 경로방향 / 경로이탈 오차
%   Delta              lookahead 거리. 수선의 발에서 **경로를 따라** 잰다
%   chi_d              요구 침로각 = pi_p - atan(y_e^p/Delta)
%   chi, psi, beta     실제 침로각, 선수각, 크랩각.  chi = psi + beta
%
% 좌표는 전부 w05-los-standard-geo.tex 에 있다. **이 파일에는 산술이 없다** —
% TikZ 의 calc 에 매크로 좌표와 음수 계수를 함께 물리면 pgf 의 계수 파서가
% 폭주하고, 이름표를 눈으로 놓으면 세 개가 시선 벡터 위에 겹친다. 설명 글은
% 그림 **아래**에 둔다 — 옆에 두면 전체 폭이 A4 본문 폭을 넘어 축소된다.
%
% 볼트의 기존 그림 w05-los-geometry.svg 는 같은 법칙을 더 작은 수치로 단순히
% 그린 것이다. 둘을 함께 싣는다 — 단순한 쪽으로 법칙을 읽고, 이 그림으로
% 표준 표기(위첨자 p, 침로각 chi)를 읽는다. 기존 그림은 건드리지 않는다.
\documentclass[border=5pt]{standalone}
% capstone-style.tex — 이 볼트의 TikZ 그림이 공유하는 스타일
%
% 쓰는 법:  % capstone-style.tex — 이 볼트의 TikZ 그림이 공유하는 스타일
%
% 쓰는 법:  % capstone-style.tex — 이 볼트의 TikZ 그림이 공유하는 스타일
%
% 쓰는 법:  \input{capstone-style}   (assets/src/ 안에서)
% 빌드:     bash _tools/tikz2svg.sh assets/src/w02-node-topic.tex assets/w02-node-topic.svg
%
% 왜 TikZ 인가
%   손으로 SVG 좌표를 쓰면 화살촉·선 굵기·글자 크기가 그림마다 미묘하게 달라진다.
%   그것이 "자료가 어설퍼 보인다"의 정체다. 스타일을 한 번만 정하고 상대좌표로
%   배치하면 좌표를 고쳐도 그림이 무너지지 않는다.
%   수식은 본문과 같은 TeX 글꼴로 나온다.
%
% 한글
%   ko.TeX(DVI 경로)를 쓴다. XeLaTeX 은 TikZ 그래픽을 PDF special 로 내보내는데
%   dvisvgm 이 그것을 받으려면 Ghostscript < 10.01 이 필요하다. 이 PC 에는
%   10.07 이 깔려 있어 그 경로가 막힌다. latex -> DVI -> dvisvgm 이 유일하게
%   동작하는 조합이다.

\usepackage[hangul]{kotex}
\usepackage{tikz}
\usepackage{amsmath}
\usepackage{xcolor}
\usetikzlibrary{arrows.meta, positioning, calc, fit, backgrounds, decorations.pathmorphing}

% ---- 색 --------------------------------------------------------------------
% 강조는 하나만. 나머지는 검정. 흑백 인쇄에서도 읽혀야 한다.
\definecolor{ink}{HTML}{1A1A1A}
\definecolor{soft}{HTML}{666666}
\definecolor{accent}{HTML}{7C3AED}
\definecolor{warn}{HTML}{B91C1C}

% 계층 색 — Simulink 모델의 블록 색과 같은 규칙을 쓴다
\definecolor{guid}{HTML}{1D4ED8}    % 유도
\definecolor{ctrl}{HTML}{EA580C}    % 제어
\definecolor{thr}{HTML}{D97706}     % 추진기
\definecolor{plant}{HTML}{16A34A}   % 운동모델
\definecolor{nav}{HTML}{7C3AED}     % 항법 · 신호처리
\definecolor{pathc}{HTML}{0D9488}   % 계획 경로 · 항적
\definecolor{errc}{HTML}{C2410C}    % 오차
\definecolor{flow}{HTML}{0369A1}    % 조류 · 바람

% ---- 선과 화살촉 -----------------------------------------------------------
\tikzset{
  sig/.style   = {ink, line width=0.5pt, -{Stealth[length=4pt, width=3pt]}},
  wire/.style  = {ink, line width=0.5pt},
  fb/.style    = {ink, line width=0.5pt, densely dashed,
                  -{Stealth[length=4pt, width=3pt]}},
  asig/.style  = {accent, line width=0.5pt, densely dashed,
                  -{Stealth[length=4pt, width=3pt]}},
  % 블록 — 흰 바탕, 직각, 검은 테두리. 색은 테두리에만 준다
  blk/.style   = {draw=ink, line width=0.55pt, fill=white,
                  minimum width=18mm, minimum height=8.5mm, inner sep=2pt, font=\small},
  gblk/.style  = {blk, draw=guid},
  cblk/.style  = {blk, draw=ctrl},
  tblk/.style  = {blk, draw=thr},
  pblk/.style  = {blk, draw=plant},
  nblk/.style  = {blk, draw=nav},
  % 합산점 · 분기점
  sum/.style   = {draw=ink, line width=0.55pt, fill=white, circle,
                  inner sep=0pt, minimum size=4.6mm, font=\scriptsize},
  dot/.style   = {ink, fill=ink, circle, inner sep=0pt, minimum size=1.5mm},
  % 글자
  sname/.style = {font=\small, inner sep=1.5pt},
  note/.style  = {font=\scriptsize, soft, inner sep=1.5pt},
  wnote/.style = {font=\scriptsize, warn, inner sep=1.5pt},
  anote/.style = {font=\scriptsize, accent, inner sep=1.5pt},
  axis/.style  = {ink, line width=0.5pt},
  tick/.style  = {font=\scriptsize, soft},
  % 묶음 상자 — 서브시스템 표시
  grp/.style   = {draw=soft, line width=0.4pt, densely dotted, rounded corners=1mm,
                  inner sep=3mm},
}
   (assets/src/ 안에서)
% 빌드:     bash _tools/tikz2svg.sh assets/src/w02-node-topic.tex assets/w02-node-topic.svg
%
% 왜 TikZ 인가
%   손으로 SVG 좌표를 쓰면 화살촉·선 굵기·글자 크기가 그림마다 미묘하게 달라진다.
%   그것이 "자료가 어설퍼 보인다"의 정체다. 스타일을 한 번만 정하고 상대좌표로
%   배치하면 좌표를 고쳐도 그림이 무너지지 않는다.
%   수식은 본문과 같은 TeX 글꼴로 나온다.
%
% 한글
%   ko.TeX(DVI 경로)를 쓴다. XeLaTeX 은 TikZ 그래픽을 PDF special 로 내보내는데
%   dvisvgm 이 그것을 받으려면 Ghostscript < 10.01 이 필요하다. 이 PC 에는
%   10.07 이 깔려 있어 그 경로가 막힌다. latex -> DVI -> dvisvgm 이 유일하게
%   동작하는 조합이다.

\usepackage[hangul]{kotex}
\usepackage{tikz}
\usepackage{amsmath}
\usepackage{xcolor}
\usetikzlibrary{arrows.meta, positioning, calc, fit, backgrounds, decorations.pathmorphing}

% ---- 색 --------------------------------------------------------------------
% 강조는 하나만. 나머지는 검정. 흑백 인쇄에서도 읽혀야 한다.
\definecolor{ink}{HTML}{1A1A1A}
\definecolor{soft}{HTML}{666666}
\definecolor{accent}{HTML}{7C3AED}
\definecolor{warn}{HTML}{B91C1C}

% 계층 색 — Simulink 모델의 블록 색과 같은 규칙을 쓴다
\definecolor{guid}{HTML}{1D4ED8}    % 유도
\definecolor{ctrl}{HTML}{EA580C}    % 제어
\definecolor{thr}{HTML}{D97706}     % 추진기
\definecolor{plant}{HTML}{16A34A}   % 운동모델
\definecolor{nav}{HTML}{7C3AED}     % 항법 · 신호처리
\definecolor{pathc}{HTML}{0D9488}   % 계획 경로 · 항적
\definecolor{errc}{HTML}{C2410C}    % 오차
\definecolor{flow}{HTML}{0369A1}    % 조류 · 바람

% ---- 선과 화살촉 -----------------------------------------------------------
\tikzset{
  sig/.style   = {ink, line width=0.5pt, -{Stealth[length=4pt, width=3pt]}},
  wire/.style  = {ink, line width=0.5pt},
  fb/.style    = {ink, line width=0.5pt, densely dashed,
                  -{Stealth[length=4pt, width=3pt]}},
  asig/.style  = {accent, line width=0.5pt, densely dashed,
                  -{Stealth[length=4pt, width=3pt]}},
  % 블록 — 흰 바탕, 직각, 검은 테두리. 색은 테두리에만 준다
  blk/.style   = {draw=ink, line width=0.55pt, fill=white,
                  minimum width=18mm, minimum height=8.5mm, inner sep=2pt, font=\small},
  gblk/.style  = {blk, draw=guid},
  cblk/.style  = {blk, draw=ctrl},
  tblk/.style  = {blk, draw=thr},
  pblk/.style  = {blk, draw=plant},
  nblk/.style  = {blk, draw=nav},
  % 합산점 · 분기점
  sum/.style   = {draw=ink, line width=0.55pt, fill=white, circle,
                  inner sep=0pt, minimum size=4.6mm, font=\scriptsize},
  dot/.style   = {ink, fill=ink, circle, inner sep=0pt, minimum size=1.5mm},
  % 글자
  sname/.style = {font=\small, inner sep=1.5pt},
  note/.style  = {font=\scriptsize, soft, inner sep=1.5pt},
  wnote/.style = {font=\scriptsize, warn, inner sep=1.5pt},
  anote/.style = {font=\scriptsize, accent, inner sep=1.5pt},
  axis/.style  = {ink, line width=0.5pt},
  tick/.style  = {font=\scriptsize, soft},
  % 묶음 상자 — 서브시스템 표시
  grp/.style   = {draw=soft, line width=0.4pt, densely dotted, rounded corners=1mm,
                  inner sep=3mm},
}
   (assets/src/ 안에서)
% 빌드:     bash _tools/tikz2svg.sh assets/src/w02-node-topic.tex assets/w02-node-topic.svg
%
% 왜 TikZ 인가
%   손으로 SVG 좌표를 쓰면 화살촉·선 굵기·글자 크기가 그림마다 미묘하게 달라진다.
%   그것이 "자료가 어설퍼 보인다"의 정체다. 스타일을 한 번만 정하고 상대좌표로
%   배치하면 좌표를 고쳐도 그림이 무너지지 않는다.
%   수식은 본문과 같은 TeX 글꼴로 나온다.
%
% 한글
%   ko.TeX(DVI 경로)를 쓴다. XeLaTeX 은 TikZ 그래픽을 PDF special 로 내보내는데
%   dvisvgm 이 그것을 받으려면 Ghostscript < 10.01 이 필요하다. 이 PC 에는
%   10.07 이 깔려 있어 그 경로가 막힌다. latex -> DVI -> dvisvgm 이 유일하게
%   동작하는 조합이다.

\usepackage[hangul]{kotex}
\usepackage{tikz}
\usepackage{amsmath}
\usepackage{xcolor}
\usetikzlibrary{arrows.meta, positioning, calc, fit, backgrounds, decorations.pathmorphing}

% ---- 색 --------------------------------------------------------------------
% 강조는 하나만. 나머지는 검정. 흑백 인쇄에서도 읽혀야 한다.
\definecolor{ink}{HTML}{1A1A1A}
\definecolor{soft}{HTML}{666666}
\definecolor{accent}{HTML}{7C3AED}
\definecolor{warn}{HTML}{B91C1C}

% 계층 색 — Simulink 모델의 블록 색과 같은 규칙을 쓴다
\definecolor{guid}{HTML}{1D4ED8}    % 유도
\definecolor{ctrl}{HTML}{EA580C}    % 제어
\definecolor{thr}{HTML}{D97706}     % 추진기
\definecolor{plant}{HTML}{16A34A}   % 운동모델
\definecolor{nav}{HTML}{7C3AED}     % 항법 · 신호처리
\definecolor{pathc}{HTML}{0D9488}   % 계획 경로 · 항적
\definecolor{errc}{HTML}{C2410C}    % 오차
\definecolor{flow}{HTML}{0369A1}    % 조류 · 바람

% ---- 선과 화살촉 -----------------------------------------------------------
\tikzset{
  sig/.style   = {ink, line width=0.5pt, -{Stealth[length=4pt, width=3pt]}},
  wire/.style  = {ink, line width=0.5pt},
  fb/.style    = {ink, line width=0.5pt, densely dashed,
                  -{Stealth[length=4pt, width=3pt]}},
  asig/.style  = {accent, line width=0.5pt, densely dashed,
                  -{Stealth[length=4pt, width=3pt]}},
  % 블록 — 흰 바탕, 직각, 검은 테두리. 색은 테두리에만 준다
  blk/.style   = {draw=ink, line width=0.55pt, fill=white,
                  minimum width=18mm, minimum height=8.5mm, inner sep=2pt, font=\small},
  gblk/.style  = {blk, draw=guid},
  cblk/.style  = {blk, draw=ctrl},
  tblk/.style  = {blk, draw=thr},
  pblk/.style  = {blk, draw=plant},
  nblk/.style  = {blk, draw=nav},
  % 합산점 · 분기점
  sum/.style   = {draw=ink, line width=0.55pt, fill=white, circle,
                  inner sep=0pt, minimum size=4.6mm, font=\scriptsize},
  dot/.style   = {ink, fill=ink, circle, inner sep=0pt, minimum size=1.5mm},
  % 글자
  sname/.style = {font=\small, inner sep=1.5pt},
  note/.style  = {font=\scriptsize, soft, inner sep=1.5pt},
  wnote/.style = {font=\scriptsize, warn, inner sep=1.5pt},
  anote/.style = {font=\scriptsize, accent, inner sep=1.5pt},
  axis/.style  = {ink, line width=0.5pt},
  tick/.style  = {font=\scriptsize, soft},
  % 묶음 상자 — 서브시스템 표시
  grp/.style   = {draw=soft, line width=0.4pt, densely dotted, rounded corners=1mm,
                  inner sep=3mm},
}

% assets/src/w05-los-standard-geo.tex — 손으로 고치지 않는다.
%
% 대학원 과목 "센서신호처리 및 융합" 의 figures/src/w05-los-geometry-geo.tex
% 를 그대로 옮긴 것이다. 그 파일은 MATLAB 스크립트가 생성한다:
%
%   matlab -batch "addpath _tools; w05_losgeo"        (대학원 볼트에서)
%
% Fossen TTK4190 의 표준 LOS 그림 배치. 좌표는 (East, North) 순이고
% 단위는 mm (1 m = 1.55 mm).

\def\WPi{(0.000,0.000)}
\def\WPnext{(118.737,99.632)}
\def\Foot{(37.996,31.882)}
\def\Ves{(53.937,12.884)}
\def\Aim{(83.116,69.742)}
\def\Cross{(53.937,45.258)}
\def\PathA{(-9.499,-7.971)}
\def\PathB{(132.985,111.588)}
\def\MidXe{(18.998,15.941)}
\def\MidYe{(45.966,22.383)}
\def\MidDel{(60.556,50.812)}
\def\MidLOS{(69.110,42.451)}
\def\RaA{(36.004,30.211)}
\def\RaB{(37.675,28.219)}
\def\RaC{(39.667,29.891)}
\def\DropA{(83.116,-13.116)}
\def\DropB{(53.937,-13.116)}
\def\RbA{(80.116,-13.116)}
\def\RbB{(80.116,-10.116)}
\def\RbC{(83.116,-10.116)}
\def\Xntop{(53.937,58.884)}
\def\RadPsi{20.00}
\def\ArcTopPsi{(53.937,32.884)}
\def\LabPsi{(56.886,36.904)}
\def\RadChi{26.00}
\def\ArcTopChi{(53.937,38.884)}
\def\LabChi{(59.699,42.530)}
\def\RadChid{33.00}
\def\ArcTopChid{(53.937,45.884)}
\def\LabChid{(62.674,49.044)}
\def\ArcTopBeta{(57.324,26.469)}
\def\LabBeta{(62.745,28.328)}
\def\ArcTopCorr{(76.758,64.407)}
\def\RadCorr{8.30}
\def\ArcMidCorr{(76.942,62.004)}
\def\HullBow{(55.388,18.706)}
\def\HullStbd{(55.559,9.883)}
\def\HullPort{(51.096,10.996)}
\def\LabXb{(60.590,39.567)}
\def\LabU{(69.296,50.899)}
\def\LabDelta{(67.525,62.665)}
\def\LabPip{(49.937,48.258)}

\def\ScrPip{40.0000}
\def\ScrChid{62.8337}
\def\ScrChi{68.0000}
\def\ScrPsi{76.0000}

\def\ValPip{50}
\def\ValChid{27.17}
\def\ValChi{22}
\def\ValPsi{14}
\def\ValBeta{8}
\def\ValXe{32}
\def\ValYe{16}
\def\ValDelta{38}
\def\ValCorr{22.83}
\def\ValLOS{41.23}
\def\Scale{1.55}


%  색은 \tikzset 보다 먼저 정의해야 한다 — 스타일이 곧바로 쓴다.
%  capstone-style 에 없는 둘만 여기서 더한다.
\definecolor{knownc}{HTML}{1D4ED8}
\definecolor{alng}{HTML}{B45309}      % 경로방향(along-track) 성분

\tikzset{
  pathline/.style = {ink, line width=0.9pt},
  aux/.style      = {soft, line width=0.35pt, densely dashed},
  errline/.style  = {errc, line width=1.0pt},
  alngline/.style = {alng, line width=1.0pt},
  losline/.style  = {accent, line width=1.0pt, -{Stealth[length=4.6pt,width=3.4pt]}},
  velline/.style  = {pathc, line width=1.0pt, -{Stealth[length=4.6pt,width=3.4pt]}},
  bodyline/.style = {ink, line width=0.8pt, -{Stealth[length=4.2pt,width=3.0pt]}},
  ang/.style      = {soft, line width=0.5pt},
  rmark/.style    = {soft, line width=0.4pt},
  %  **식을 그림 안에 크게** 둔다. 상자에 넣어야 그림의 선들과 섞이지 않는다.
  %  색은 그 식이 무엇인지에 따라 다르다:
  %    known  이미 아는 값        knownc 파랑
  %    want   구하려는 값         accent 보라
  %    limit  극한에서 되는 것    pathc  청록
  known/.style = {draw=knownc, line width=0.6pt, fill=white, inner sep=3.2pt,
                  rounded corners=1pt, font=\small},
  want/.style  = {draw=accent, line width=0.9pt, fill=white, inner sep=3.4pt,
                  rounded corners=1pt, font=\small},
  limit/.style = {draw=pathc, line width=0.6pt, fill=white, inner sep=3.2pt,
                  rounded corners=1pt, font=\small},
  capt/.style  = {font=\scriptsize\bfseries, warn},
  lead/.style  = {soft, line width=0.35pt, -{Stealth[length=3.4pt,width=2.4pt]}},
}

\begin{document}
\begin{tikzpicture}[x=1mm, y=1mm]

% ---- 북·동 축 ------------------------------------------------------------
\draw[axis, -{Stealth[length=4pt,width=3pt]}] (2,96) -- (2,109);
\node[font=\scriptsize, anchor=south] at (2,109.5) {$x^n$~(북)};
\draw[axis, -{Stealth[length=4pt,width=3pt]}] (2,96) -- (15,96);
\node[font=\scriptsize, anchor=west] at (15.5,96) {$y^n$~(동)};

% ---- 경로 ----------------------------------------------------------------
\draw[pathline] \PathA -- \PathB;
\fill[ink] \WPi    circle (1.5pt);
\fill[ink] \WPnext circle (1.5pt);
\node[font=\small, anchor=east]        at (-3.0,1.0)    {$\mathbf{p}_{i}^{n}$};
%  p_i 의 오른쪽 아래에 둔다. 왼쪽에 두면 그림 왼쪽 끝에 딱 붙어 잘려 보인다.
\node[note, anchor=north west]         at (2.5,-2.0)    {$[\,x_i^n,\ y_i^n\,]^{\!\top}$};
\node[font=\small, anchor=south east]  at (117.0,104.5) {$\mathbf{p}_{i+1}^{n}$};
%  두 웨이포인트 모두 성분을 적는다. 하나만 적으면 다른 하나가 무엇인지
%  독자가 유추해야 한다.
%  경로선 **아래** 오른쪽에 둔다. 선 위에 두면 글자가 선에 물린다.
\node[note, anchor=north west]         at (119.5,96.0)  {$[\,x_{i+1}^n,\ y_{i+1}^n\,]^{\!\top}$};

% ---- 선박을 지나는 북쪽 점선 (psi, chi, chi_d 가 여기서 잰다) -----------
\draw[aux] \Ves -- \Xntop;
\node[font=\small, soft, anchor=south east] at (53.0,59.0) {$x^n$};

% ---- pi_p 는 **웨이포인트에서** 잰다 ------------------------------------
%  pi_p 는 다리 자체의 방위각이므로 다리가 시작하는 p_i 에서 북에서부터 재야
%  읽힌다. 반지름 14 는 x_e^p 이름표(17,17.8)를 피한 값이다.
\draw[aux] \WPi -- (0,22);
\node[font=\small, soft, anchor=south] at (0,22.5) {$x^n$};
\draw[ang] (0,14) arc (90:\ScrPip:14);
\node[font=\small] at (4.2,9.0) {$\pi_p$};

% ---- 수선의 발, 경로방향 오차와 경로이탈 오차 ---------------------------
\draw[alngline] \WPi -- \Foot;
\draw[errline]  \Foot -- \Ves;
\fill[ink] \Foot circle (1.1pt);
\draw[rmark] \RaA -- \RaB -- \RaC;

\node[alng, font=\small, anchor=south east] at (17.0,17.8) {$x_e^{p}$};
\node[errc, font=\small, anchor=east]       at (43.6,22.4) {$y_e^{p}$};

% ---- lookahead 와 조준점 -------------------------------------------------
\draw[alngline] \Foot -- \Aim;
\fill[accent] \Aim circle (1.4pt);
\node[alng, font=\small, anchor=south east] at \LabDelta {$\Delta$};
%  Delta 가 무엇인지 말로도 짚어 준다. 화살표는 경로 위의 점 (66,55.4) 로
%  보낸다 — 66 * tan(40도) = 55.4 이므로 그 점은 수선의 발–조준점 선분 위에 있다.
\node[capt,anchor=east] at (55.0,66.0) {lookahead 거리};
\draw[lead] (56.0,65.5) -- (66.0,55.4);
%  Delta 가 **어디서 어디까지인지**를 양 끝에 적는다. 그래야 x_e^p 와 섞이지 않는다.
\node[note, anchor=south east] at (36.5,33.5) {수선의 발};
\node[note, anchor=south west] at (84.0,73.5) {조준점};

% ---- 시선 벡터 -----------------------------------------------------------
\draw[losline] \Ves -- \Aim;
\node[accent, font=\scriptsize, anchor=north west, rotate=\ScrChid]
     at (70.8,41.0) {시선 벡터};

%  조준점에서의 각. 두 선을 조준점에서 **뒤로** 본 구간이다.
\draw[ang] \ArcTopCorr arc (\ScrPip+180:\ScrChid+180:\RadCorr);
\draw[soft, line width=0.35pt, -{Stealth[length=3.4pt,width=2.4pt]}]
     (98.0,52.0) -- \ArcMidCorr;
\node[font=\scriptsize, anchor=west] at (98.5,51.6)
     {$\tan^{-1}\!\big(y_e^{p}/\Delta\big)$};

%  조준점에서 내린 세로 점선, 선박에서 오는 가로 점선, 만나는 곳의 직각
\draw[aux] \Aim   -- \DropA;
\draw[aux] \Ves   -- \DropB;
\draw[rmark] \RbA -- \RbB -- \RbC;

% ---- 선박 ----------------------------------------------------------------
\fill[ink, opacity=0.85] \HullBow -- \HullStbd -- \HullPort -- cycle;

\draw[bodyline] \Ves -- ++(\ScrPsi:24);
\node[font=\small] at \LabXb {$x_b$};

\draw[velline] \Ves -- ++(\ScrChi:38);
\node[pathc, font=\small] at \LabU {$U$};

\draw[ang] \ArcTopPsi  arc (90:\ScrPsi:\RadPsi);
\node[font=\small]       at \LabPsi  {$\psi$};
\draw[ang] \ArcTopChi  arc (90:\ScrChi:\RadChi);
\node[pathc, font=\small] at \LabChi {$\chi$};
\draw[ang] \ArcTopChid arc (90:\ScrChid:\RadChid);
\node[accent, font=\small] at \LabChid {$\chi_d$};

\draw[accent, line width=0.7pt] \ArcTopBeta arc (\ScrPsi:\ScrChi:14);
\node[anote] at \LabBeta {$\beta$};

\node[note, anchor=north east] at (51.0,8.0) {$(x^n,\ y^n)$};

% ===================== 그림 안의 식 ======================================
%  아는 값은 파란 상자, 구하려는 값은 보라 상자, 극한에서 성립하는 것은 청록
%  상자. 상자마다 무엇인지 한 낱말로 적는다.

%  경로 접선각 — 두 웨이포인트에서 바로 나온다
\node[known, anchor=north west] at (-6,-32)
     {$\pi_p = \operatorname{atan2}\!\big(y_{i+1}^n - y_i^n,\ x_{i+1}^n - x_i^n\big)$};
\node[capt,anchor=north west] at (-6,-40.5) {아는 값};

%  경로이탈 오차 — 위치와 경로에서 바로 나온다
\node[known, anchor=north west] at (-6,-48)
     {$y_e^{p} = -\big(x^n - x_i^n\big)\sin\pi_p + \big(y^n - y_i^n\big)\cos\pi_p$};
\node[capt,anchor=north west] at (-6,-56.5) {아는 값};

%  웨이포인트 둘
\node[known, anchor=north west] at (-6,-64)
     {$\mathbf{p}_i^{n} = [\,x_i^n,\ y_i^n\,]^{\!\top}$,\quad
       $\mathbf{p}_{i+1}^{n} = [\,x_{i+1}^n,\ y_{i+1}^n\,]^{\!\top}$};
\node[capt,anchor=north west] at (-6,-72.5) {아는 값};

%  구하려는 것
\node[want, anchor=north west] at (76,-32)
     {$\chi_d = \pi_p - \tan^{-1}\!\big(K_p\,y_e^{p}\big),
       \qquad K_p = \dfrac{1}{\Delta}$};
\node[capt,anchor=north west] at (76,-46) {구하려는 값};

%  극한에서 성립하는 것
\node[limit, anchor=north west] at (76,-56)
     {$y_e^{p} \to 0, \qquad \chi_d \to \pi_p$};
\node[anchor=north west, align=left, text width=64mm, font=\scriptsize, soft] at (76,-65)
     {경로 위에서는 보정이 사라지고 선박이 구간을 따라 곧게 나아감.
      이 법칙이 노리는 바를 극한으로 적은 것임.};

% ===================== 아래 : 법칙과 수치 ================================
\draw[soft, line width=0.3pt] (-6,-80) -- (146,-80);

\node[anchor=north west, align=left, text width=68mm, font=\scriptsize] at (-6,-84) {%
  \textbf{항은 둘뿐임.} $\pi_p$ 는 선박을 경로와 \emph{나란히} 세우고,
  $-\tan^{-1}(y_e^{p}/\Delta)$ 는 선박을 경로 \emph{쪽으로} 기울임.
  그 밖에 들어가는 항이 없으므로 이 법칙이 못 하는 일도 분명함 (1-7절).\\[4pt]
  \textbf{보정각은 $90^\circ$ 를 넘지 못함.} $y_e^{p}\to 0$ 이면 선박이 경로
  위에 있고 $\chi_d\to\pi_p$ 임. $y_e^{p}\to\infty$ 이면
  $\chi_d\to\pi_p-90^\circ$ 가 되어 경로를 정면으로 향함 — 경로에서 멀어지는
  지령은 나오지 않음.\\[4pt]
  \textbf{$\Delta$ 는 경로를 따라 잼.} 선박이 아니라 수선의 발에서 재는
  거리임. 선박에서 조준점까지의 실제 거리는 $\sqrt{\Delta^2+(y_e^{p})^2}$
  이고 $\Delta$ 보다 큼.};

\node[anchor=north west, align=left, text width=68mm, font=\scriptsize] at (76,-84) {%
  {\small\bfseries 표기와 이 그림의 수치}\\[4pt]
  웨이포인트는 $\mathbf{p}_i^{n} = [\,x_i^n,\ y_i^n\,]^{\!\top}$ 이고 $\{n\}$
  에서 읽은 값임. 오차 둘에 위첨자 $p$ 가 붙는 까닭은 그 둘이 \textbf{경로}
  좌표계 $\{p\}$ 에서 읽은 값이기 때문임 — $\{n\}$ 도 $\{b\}$ 도 아님.
  위첨자가 세 좌표계를 갈라 놓음.\\[4pt]
  그림은 $\pi_p = \ValPip^\circ$, $x_e^{p} = \ValXe$\,m,
  $y_e^{p} = \ValYe$\,m, $\Delta = \ValDelta$\,m 로 그린 것임. 보정각은
  $\tan^{-1}(\ValYe/\ValDelta) = \ValCorr^\circ$ 이고
  $\chi_d = \ValChid^\circ$ 임. 시선 거리는 $\ValLOS$\,m 임.
  선체는 $\psi = \ValPsi^\circ$ 에 크랩각 $\beta = \ValBeta^\circ$ 이므로
  $\chi = \ValChi^\circ$ 임 — 아직 지령 침로에 올라서지 못한 상태이고, 그래서
  선체에 걸린 호 셋이 서로 다름.\\[4pt]
  \textbf{침로각인가 선수각인가.} 적힌 그대로의 법칙은 \emph{침로각} $\chi_d$
  를 지령함. 이 침로각을 받는 것은 4주차의 선수각 제어기이므로, 실습에서는
  $\chi_d$ 자리에 $\psi_d$ 를 넣은 같은 식을 씀. 둘은 $\beta = 0$ 일 때만
  정확히 같음. 그 차이는 1-2절이 정의하고, 1-7절이 경로에서 치르는 값을 잼.};

\end{tikzpicture}
\end{document}
